% Environment, GitOps pattern, use case

\section{Implementation}
\label{sec:implementation}

The platform runs on k3s on a single node \cite{kubernetes2024}. The node count is an implementation choice, and moving to a multi-node cluster should change only the Kustomize overlay, a move not yet performed. Library versions and base images are pinned to the highest stable release as of April 2026 so that testing is reproducible, and a component is intended to be replaceable by an equivalent without changing the design.

Components are grouped one namespace per layer and joined by stable API contracts, which lets network policy default to deny with declared exceptions. Argo CD is the only path by which change reaches the cluster \cite{limoncelli2022gitops}, with every component rendered from an ApplicationSet template so adding a component adds no boilerplate. Tekton builds and tests. Istio enforces mutual TLS in strict mode on internal traffic, Kyverno validates at admission, and Falco watches the kernel at runtime. The default is one idle replica per workload, with autoscaling templates prepared but not exercised under load.

Platform code stays domain-agnostic in its own directory, while the domain lives in a separate use-case directory applied through a Kustomize overlay setting only the name prefix, registry mapping, initial replicas, and autoscaling bounds.

The verification use case is crypto asset market analytics on Coinbase data, chosen for three reasons. It exercises all nine layers at once. Crypto markets trade continuously with no closing hours \cite{makarov2020trading}, so the ingestion pipeline and stream processing are driven by real event flow around the clock rather than synthetic load on a schedule. Price volatility is high and clustered \cite{katsiampa2017volatility}, so feature distributions are expected to shift within hours, which makes the domain suitable for drift work. The drift result in Section~\ref{subsec:results} nevertheless comes from a controlled injection, not from natural shift. The use case runs two ingestion pipelines. A WebSocket collector feeds Flink for windowed features, and a REST collector backfills one-minute candles from 1 March 2026.
